\labsection{9}{Anomaly detection and limited-label learning}{sec:lab09-anomaly-labels}

\begin{labproject}{Learn when labels are absent, rare, or expensive}
\textbf{Coverage.} Chapters 18--19.\\
\textbf{Core data.} Credit-card transactions and Breast Cancer Wisconsin.\\
\textbf{Goal.} Evaluate an unsupervised anomaly ranking with labels held back from fitting, then simulate self-training and active learning under a fixed labeling budget.

\begin{labmode}
The common theme is label scarcity. In the anomaly activity, labels are unavailable to the fit but available for evaluation; in active learning, labels can be purchased selectively and must not be replaced by repeated test-set inspection.
\end{labmode}

\medskip\noindent\textbf{Part A --- Isolation Forest anomaly ranking.}\par
\begin{lstlisting}[style=mlpython]
from pathlib import Path
import pandas as pd
import matplotlib.pyplot as plt
from sklearn.ensemble import IsolationForest
from sklearn.metrics import average_precision_score, roc_auc_score
from sklearn.model_selection import train_test_split
\end{lstlisting}

\begin{lstlisting}[style=mlpython]
PROJECT_ROOT = Path.cwd().parent
train_path = PROJECT_ROOT / "all_datasets" / "credit_card_fraud_dataset" / "credit_card_fraud_10k.csv"
df = pd.read_csv(train_path)
\end{lstlisting}

\begin{lstlisting}[style=mlpython]
# Use only behavior/risk features. The transaction ID is an identifier, not a signal.
feature_cols = [
    "amount",
    "transaction_hour",
    "foreign_transaction",
    "location_mismatch",
    "device_trust_score",
    "velocity_last_24h",
    "cardholder_age",
]
X = df[feature_cols]
y = df["is_fraud"]

X_train, X_test, y_train, y_test = train_test_split(
    X, y, test_size=0.30, random_state=42, stratify=y
)

# The fraud labels are deliberately NOT passed to fit.
model = IsolationForest(
    n_estimators=300, contamination="auto", random_state=42, n_jobs=-1
)
model.fit(X_train)
anomaly_score = -model.decision_function(X_test)

prevalence = y_test.mean()
ap = average_precision_score(y_test, anomaly_score)
print("fraud prevalence:", round(prevalence, 4))
print("ROC AUC:", round(roc_auc_score(y_test, anomaly_score), 4))
print("average precision:", round(ap, 4))
print("AP / prevalence:", round(ap / prevalence, 2))
\end{lstlisting}

\begin{lstlisting}[style=mlpython]
ranked = X_test.copy()
ranked["fraud_label"] = y_test.to_numpy()
ranked["anomaly_score"] = anomaly_score
ranked = ranked.sort_values("anomaly_score", ascending=False)
review_budget = min(100, len(ranked))
top_k_precision = ranked.head(review_budget)["fraud_label"].mean()
print(f"precision among top {review_budget}:", round(top_k_precision, 4))
print(ranked[["fraud_label", "anomaly_score", "amount"]].head(20))
\end{lstlisting}

\begin{lstlisting}[style=mlpython]
plt.figure(figsize=(6, 4))
plt.hist(anomaly_score[y_test.to_numpy() == 0], bins=30, alpha=0.65, label="non-fraud")
plt.hist(anomaly_score[y_test.to_numpy() == 1], bins=30, alpha=0.65, label="fraud")
plt.xlabel("Anomaly score")
plt.ylabel("Transactions")
plt.title("Isolation Forest score distribution")
plt.legend()
plt.tight_layout()
plt.show()
\end{lstlisting}

\medskip\noindent\textbf{Part B --- Self-training and active learning.}\par
\begin{lstlisting}[style=mlpython]
from pathlib import Path
import numpy as np
import pandas as pd
import matplotlib.pyplot as plt
from sklearn.linear_model import LogisticRegression
from sklearn.metrics import accuracy_score
from sklearn.model_selection import train_test_split
from sklearn.pipeline import Pipeline
from sklearn.preprocessing import StandardScaler
from sklearn.semi_supervised import SelfTrainingClassifier
\end{lstlisting}

\begin{lstlisting}[style=mlpython]
rng = np.random.default_rng(42)
rng_random = np.random.default_rng(123)

PROJECT_ROOT = Path.cwd().parent
train_path = PROJECT_ROOT / "all_datasets" / "breast_cancer_dataset" / "data.csv"
df = pd.read_csv(train_path)

X = df.drop(columns=["id", "diagnosis", "Unnamed: 32"], errors="ignore")
y = (df["diagnosis"] == "M").astype(int).to_numpy()
\end{lstlisting}

\begin{lstlisting}[style=mlpython]
# The outer test set stays untouched until the label-acquisition policy is finished.
X_dev, X_test, y_dev, y_test = train_test_split(
    X, y, test_size=0.25, random_state=42, stratify=y
)
X_pool, X_val, y_pool, y_val = train_test_split(
    X_dev, y_dev, test_size=0.25, random_state=7, stratify=y_dev
)
X_pool = X_pool.reset_index(drop=True)
y_pool = np.asarray(y_pool)
\end{lstlisting}

\begin{lstlisting}[style=mlpython]
# ----- Semi-supervised self-training -----
# Pretend only 30% of the pool labels are initially available.
y_partial = y_pool.copy()
hide = rng.random(len(y_partial)) < 0.70
y_partial[hide] = -1

semi = Pipeline([
    ("scale", StandardScaler()),
    ("model", SelfTrainingClassifier(
        LogisticRegression(max_iter=2000), threshold=0.90
    )),
])
semi.fit(X_pool, y_partial)
print("self-training validation accuracy:",
      round(accuracy_score(y_val, semi.predict(X_val)), 3))
\end{lstlisting}

\begin{lstlisting}[style=mlpython]
# ----- Active-learning simulation -----
# Start from the same small class-balanced labeled set for both policies.
class0 = np.where(y_pool == 0)[0]
class1 = np.where(y_pool == 1)[0]
initial = np.concatenate([
    rng.choice(class0, size=20, replace=False),
    rng.choice(class1, size=20, replace=False),
])
active_labeled = set(initial.tolist())
random_labeled = set(initial.tolist())

label_counts = []
active_val_scores = []
random_val_scores = []
query_batch = 15
rounds = 8

def make_model():
    return Pipeline([
        ("scale", StandardScaler()),
        ("model", LogisticRegression(max_iter=2000)),
    ])
\end{lstlisting}

\begin{lstlisting}[style=mlpython]
for round_no in range(rounds):
    active_idx = np.array(sorted(active_labeled))
    random_idx = np.array(sorted(random_labeled))

    active = make_model()
    random_model = make_model()
    active.fit(X_pool.iloc[active_idx], y_pool[active_idx])
    random_model.fit(X_pool.iloc[random_idx], y_pool[random_idx])

    active_score = accuracy_score(y_val, active.predict(X_val))
    random_score = accuracy_score(y_val, random_model.predict(X_val))
    label_counts.append(len(active_labeled))
    active_val_scores.append(active_score)
    random_val_scores.append(random_score)
    print(
        f"round {round_no + 1}: labels={len(active_labeled)}, "
        f"active_val={active_score:.3f}, random_val={random_score:.3f}"
    )

    if round_no == rounds - 1:
        break

    # Uncertainty sampling: ask for labels closest to probability 0.5.
    active_unlabeled = np.array([
        i for i in range(len(X_pool)) if i not in active_labeled
    ])
    prob = active.predict_proba(X_pool.iloc[active_unlabeled])[:, 1]
    uncertainty = np.abs(prob - 0.5)
    active_query = active_unlabeled[np.argsort(uncertainty)[:query_batch]]
    active_labeled.update(active_query.tolist())

    # Same label budget, but query randomly as a baseline policy.
    random_unlabeled = np.array([
        i for i in range(len(X_pool)) if i not in random_labeled
    ])
    random_query = rng_random.choice(
        random_unlabeled, size=query_batch, replace=False
    )
    random_labeled.update(random_query.tolist())
\end{lstlisting}

\begin{lstlisting}[style=mlpython]
plt.figure(figsize=(6, 4))
plt.plot(label_counts, active_val_scores, marker="o", label="uncertainty sampling")
plt.plot(label_counts, random_val_scores, marker="o", label="random querying")
plt.xlabel("Number of labeled pool examples")
plt.ylabel("Validation accuracy")
plt.title("Label efficiency during development")
plt.legend()
plt.tight_layout()
plt.show()
\end{lstlisting}

\begin{lstlisting}[style=mlpython]
# The acquisition policy is now locked. Evaluate each final model on test once.
active_final_idx = np.array(sorted(active_labeled))
random_final_idx = np.array(sorted(random_labeled))
active_final = make_model()
random_final = make_model()
active_final.fit(X_pool.iloc[active_final_idx], y_pool[active_final_idx])
random_final.fit(X_pool.iloc[random_final_idx], y_pool[random_final_idx])
print("final active-learning test accuracy:",
      round(accuracy_score(y_test, active_final.predict(X_test)), 3))
print("final random-query test accuracy:",
      round(accuracy_score(y_test, random_final.predict(X_test)), 3))
\end{lstlisting}

\end{labproject}

\begin{labdeliverable}
Submit fraud prevalence, anomaly-ranking metrics and a top-$k$ review result, then a label-efficiency curve comparing uncertainty sampling with an equal-budget random-query baseline. Finish with one untouched-test evaluation after the acquisition policy is locked.
\end{labdeliverable}
